Um die Lorentz-Invarianz des Maßes 

\[
\frac{d^3 k}{(2 \pi)^3 2 E(\mathbf{k})}
\]

zu zeigen, beginnen wir mit dem Lorentz-invarianten Ausdruck 

\[
\frac{d^4 k}{(2 \pi)^3} \delta\left(k^2 - m^2\right) \theta\left(k_0\right).
\]

Hierbei ist $d^4 k = dk_0 \, d^3 k$ das vierdimensionale Volumenelement im Impulsraum. Die Dirac-Delta-Funktion $\delta\left(k^2 - m^2\right)$ stellt sicher, dass wir nur über die Massenschale integrieren, und $\theta\left(k_0\right)$ berücksichtigt nur positive Energien. Die Delta-Funktion kann umgeschrieben werden als 

\[
\delta\left(k^2 - m^2\right) = \delta\left(k_0^2 - \mathbf{k}^2 - m^2\right).
\]

Verwenden wir die Identität für die Delta-Funktion:

\[
\delta\left(x^2 - x_0^2\right) = \frac{1}{2|x|}\left[\delta\left(x - x_0\right) + \delta\left(x + x_0\right)\right],
\]

so erhalten wir

\[
\delta\left(k_0^2 - \mathbf{k}^2 - m^2\right) = \frac{1}{2|k_0|} \left[\delta\left(k_0 - E(\mathbf{k})\right) + \delta\left(k_0 + E(\mathbf{k})\right)\right].
\]

Da $\theta(k_0)$ nur positive Energien berücksichtigt, bleibt nur der Term $\delta\left(k_0 - E(\mathbf{k})\right)$ übrig. Daher haben wir

\[
\delta\left(k^2 - m^2\right) \theta\left(k_0\right) = \frac{1}{2E(\mathbf{k})} \delta\left(k_0 - E(\mathbf{k})\right).
\]

Setzen wir dies in das ursprüngliche Maß ein, erhalten wir

\[
\frac{d^4 k}{(2 \pi)^3} \delta\left(k^2 - m^2\right) \theta\left(k_0\right) = \frac{dk_0 \, d^3 k}{(2 \pi)^3} \frac{1}{2E(\mathbf{k})} \delta\left(k_0 - E(\mathbf{k})\right).
\]

Integration über $k_0$ führt zu

\[
\frac{d^3 k}{(2 \pi)^3 2E(\mathbf{k})},
\]

was zeigt, dass dieses Maß Lorentz-invariant ist.

Um die Invarianz von $2 E(\mathbf{k}) \delta^3\left(\mathbf{k} - \mathbf{k}^{\prime}\right)$ zu zeigen, beachten wir, dass die Delta-Funktion $\delta^3\left(\mathbf{k} - \mathbf{k}^{\prime}\right)$ unter Lorentz-Transformationen invariant ist, da sie nur von den räumlichen Komponenten abhängt. Die Invarianz von $E(\mathbf{k})$ folgt aus der Energie-Masse-Beziehung $E(\mathbf{k}) = \sqrt{\mathbf{k}^2 + m^2}$, die unter Lorentz-Transformationen invariant bleibt, da sowohl $\mathbf{k}^2$ als auch $m^2$ invariant sind. Daher bleibt auch das Produkt $2 E(\mathbf{k}) \delta^3\left(\mathbf{k} - \mathbf{k}^{\prime}\right)$ invariant.

%CHECK: Adressiert die fehlende Herleitung der Invarianz von $2 E(\mathbf{k})$ und die vollständige Darstellung der Delta-Funktion unter Lorentz-Transformationen.